\chapter{إحاطة الأمة بلغة العرب وسنة النبي}

\section{سعة لغة العرب وحفظ السنة}
لسان العرب هو أوسع الألسنة مذهباً وأكثرها ألفاظاً، ولا يوجد رجل أحاط بلغة العرب إحاطة تامة إلا أن يكون نبياً. وكذلك لا يوجد عالم أحاط بسنة النبي صلى الله عليه وسلم إحاطة تامة بمفرده. ولكن كما أن علم اللسان لا يغيب عن مجموع العرب والعلماء، فإن علم السنة لا يغيب عن مجموع العلماء قاطبة.

وهذا مقصد عظيم جداً؛ أن يعلم العلماء أن السنة والعلم لا يضيعان في مجموعهم، بشرط أن يتقوا الله سبحانه وتعالى وأن يكونوا على حرص على الخير. 

\section{عصمة الأمة والاجتهاد الجماعي}
من المشاهد أن إحاطة الفرد الواحد بالسنة من الصعوبة بمكان، فالعالم قد ينسى وقد يجهل بعض السنن، كما حدث لأمير المؤمنين عمر بن الخطاب رضي الله عنه، وهو من هو في العلم والفضل، حين خفي عليه حديث الاستئذان ثلاثاً، ونسي مسألة التيمم للمجنب. فالعالم الفرد ينسى ويجهل، لكن مجموع الأمة معصومة عن الخطأ والضلال، ومنهج السلف الصالح معصوم عن الخطأ والضلال. 

والمقصد العظيم الذي يهدف إليه الشافعي رحمه الله تعالى هو أن الاجتهاد يكون جماعياً، بشرط أن يكون العلماء أتقياء أوفياء لدينهم، لا تدخلهم عصبية، ولا يدخلهم تقليد، ولا اتباع شهوة ولا هوى. بل يكون قصدهم نصرة سنة النبي صلى الله عليه وسلم، متبعين للشرع غير متأثرين بواقع منحرف، ولا خائفين من أحد إلا من الله عز وجل.

\section{كيف يحيط طالب العلم بالسنة؟}
من أراد أن يكاد يحيط علماً بالسنة، فعليه بالكتب الستة (صحيح البخاري، وصحيح مسلم، وسنن أبي داود، وسنن النسائي، وسنن الترمذي، وسنن ابن ماجه). ويضم إليها: مسند الإمام أحمد، ومسند أبي يعلى، ومسند البزار، ومعاجم الطبراني، وسنن الدارمي، وصحيحي ابن خزيمة وابن حبان، ومستدرك الحاكم، وسنن البيهقي.

من أحاط بهذه الدواوين يكاد أن يحيط بالسنة، وهيهات ألا ينسى أو يجهل أو يغفل، لكنه في النهاية يكون قد وصل إلى علم راسخ في سنة رسول الله صلى الله عليه وسلم. فرحم الله الشافعي وطيب ثراه، ونسأل الله أن يؤدبنا بأدبه ويفقهنا بفقهه ويرزقنا الإخلاص في القول والعمل.

\section{وصية للعوام بالرجوع للعلماء والصدق مع الله}
أما الكلام السابق فهو موجه لطالب العلم، أما العامي فنقول له: الزم كبار العلماء، فإذا ألمت بك ملمة فاسأل الكبار. 
ومع ذلك، فهذا العامي لو صدق مع الله لكان خيراً له، فالله عز وجل يقول: ﴿وَلَقَدْ يَسَّرْنَا الْقُرْآنَ لِلذِّكْرِ فَهَلْ مِن مُّدَّكِرٍ﴾. ولو أنه استفتى قلبه بصدق، كما قال النبي صلى الله عليه وسلم: «استفت قلبك وإن أفتاك المفتون وأفتوك»، وكان وقوفه وسؤاله بصدق كأنه واقف بين يدي الله عز وجل؛ هداه الله، كما وعد ربنا: ﴿وَالَّذِينَ جَاهَدُوا فِينَا لَنَهْدِيَنَّهُمْ سُبُلَنَا﴾.

\section{خاتمة ودعاء}
اللهم اقسم لنا من خشيتك ما تحول به بيننا وبين معاصيك، ومن طاعتك ما تبلغنا بها جنتك، ومن اليقين ما تهون به علينا مصائب الدنيا. اللهم متعنا بأسماعنا وأبصارنا وقوتنا ما أحييتنا، واجعله الوارث منا، واجعل ثأرنا على من ظلمنا، وانصرنا على من عادانا، اللهم لا تجعل مصيبتنا في ديننا، ولا تجعل الدنيا أكبر همنا ولا مبلغ علمنا، ولا النار مصيرنا يا أرحم الراحمين.

وصلى الله وسلم وبارك على سيد الأولين والآخرين وعلى آله وصحبه وسلم.